El backend se descompone en trece aplicaciones Django, cada una con un
dominio funcional propio, sus modelos, sus servicios y sus pruebas. Esta
división facilita el mantenimiento (cada cambio se circunscribe a una
aplicación), hace explícitos los acoplamientos entre dominios (a través
de \textit{imports} y claves foráneas) y permite el desarrollo paralelo
sobre aplicaciones independientes. La \cref{table:des-apps} resume la % TODO cuadro en vez de tabla
responsabilidad de cada una.

\begin{table}[Aplicaciones Django del backend]%
            {table:des-apps}%
            {Aplicaciones Django del backend del SCDEE.}
  \small
  \renewcommand{\arraystretch}{1.3}
  \begin{tabular}{p{2.5cm} p{11cm}}
    \hline
    \textbf{Aplicación} & \textbf{Responsabilidad} \\
    \hline
    \texttt{core} &
      Utilidades transversales: modelos abstractos
      (\texttt{UUIDModel}, \texttt{TimestampedModel},
      \texttt{OrganizationOwnedModel}), \textit{middleware}
      (contexto multi-tenant, anti-caché), gestor de contexto del
      \dfn{tenant}, paginación, \textit{throttling},
      \textit{hooks} de OpenAPI, \textit{health checks}. \\
    \texttt{accounts} &
      Autenticación de usuarios, modelo \texttt{User},
      \dfnpl{plugin} de autenticación intercambiables
      (\texttt{jwtV1}, configuraciones preparatorias para
      \acs{saml} y \acs{oidc}), \textit{backend} \acs{drf},
      lista negra de \textit{tokens} revocados. \\
    \texttt{organizations} &
      Modelo \texttt{Organization} y su configuración por
      \dfn{tenant} (\texttt{OrganizationConfig}); búsqueda y
      administración de organizaciones por el superadministrador. \\
    \texttt{subjects} &
      Asignaturas (\texttt{Subject}), grupos
      (\texttt{SubjectGroup}) y membresías con rol
      (\texttt{SubjectMembership}); permisos contextuales por
      asignatura. \\
    \texttt{courses} &
      Curso académico (\texttt{AcademicCourse}) como contenedor
      temporal de las asignaturas. \\
    \texttt{exams} &
      Definición de exámenes: \texttt{Exam}, \texttt{ExamModel},
      \texttt{PageProfile}, \texttt{RecognitionZone},
      \texttt{Problem}, \texttt{RubricCriterion},
      \texttt{ExamConvocation}. \\
    \texttt{instances} &
      Ejecución y ciclo de vida de exámenes: \texttt{ExamInstance},
      \texttt{ExamPage}, máquinas de estados de instancia y
      página. \\
    \texttt{ingestion} &
      Recepción de páginas escaneadas: reconocedores como
      \dfnpl{plugin}, \textit{dispatcher} de asignación de páginas
      a \texttt{ExamInstance}, comando de gestión
      \texttt{sftp\_watcher}, tareas Celery del \textit{pipeline}
      de reconocimiento. \\
    \texttt{annotations} &
      Anotaciones de corrección sobre las páginas escaneadas
      (\texttt{Annotation}, \texttt{AnnotationType}): audio, texto,
      formas geométricas. \\
    \texttt{grading} &
      Reglas de asignación corrector-instancia
      (\texttt{AssignmentRule}) y registro de calificaciones
      (\texttt{Grade}). \\
    \texttt{reviews} &
      Revisión por parte del alumno de exámenes corregidos:
      \texttt{ExamReview}, \texttt{ReviewRequest}, ventanas de
      revisión. \\
    \texttt{notifications} &
      Notificaciones por correo electrónico, plantillas
      internacionalizadas, tareas Celery de envío. \\
    \texttt{audit} &
      Log de auditoría de eventos del sistema
      (\texttt{AuditLog}), accesible al superadministrador. \\
    \hline
  \end{tabular}
\end{table}

El listado completo de servicios expuestos por cada aplicación, con la
firma y el comportamiento de cada uno, se incluye en
\cref{AP:SERVICES}. En el cuerpo de la memoria se desarrollan únicamente
los flujos críticos del sistema (ingesta y reconocimiento en
\cref{SEC:DES-INGEST}, ciclo de vida de la instancia en
\cref{SEC:DES-LIFECYCLE}) y los aspectos de seguridad y de contrato
\acs{api} en \cref{SEC:DES-SECURITY,SEC:DES-API}.
